%% Journal of Systems and Software regular research article
%% Stage 1 structural skeleton only; no substantive manuscript prose.

\documentclass[a4paper,fleqn]{cas-dc}

\usepackage[authoryear,longnamesfirst]{natbib}

% British English is required throughout this manuscript.
% No ORCID identifiers were supplied in the authoritative previous manuscript.
\let\printorcid\relax

% Future figure constraints:
% - PDF/vector artwork is preferred.
% - Use figure or figure*; number and cite figures in order.
% - Retain separate artwork under figures/.
%
% Future table constraints:
% - Use editable LaTeX tables with booktabs.
% - Do not use vertical rules or shading.
% - Use table or table* as appropriate.

\begin{document}

\shorttitle{Contextual Input and Shared Representation}
\shortauthors{A. Kadouma et al.}

\title[mode=title]{Evaluating Contextual Input and Shared Representation in Multi-Metric Runtime Forecasting for Cloud Software Systems}

\author[1]{Abdelhak Kadouma}
\cormark[1]
\ead{abdelhak.kadouma@uwasa.fi}
\credit{Conceptualization, Methodology, Software, Data curation, Formal analysis, Validation, Visualization, Writing -- original draft}

\author[1]{Amit Shukla}
\ead{amit.shukla@uwasa.fi}
\credit{Supervision, Conceptualization, Methodology, Writing -- review \& editing}

\author[1]{Mohammed Elmusrati}
\ead{moel@uwasa.fi}
\credit{Supervision, Project administration, Writing -- review \& editing}

\affiliation[1]{
  organization={University of Vaasa},
  city={Vaasa},
  country={Finland}
}

\cortext[1]{Corresponding author}

\begin{abstract}
Software engineers designing multi-metric runtime forecasters must decide
whether to include other resources' telemetry and whether multiple targets
should share a predictive representation. These decisions affect forecasting
accuracy and model consolidation. Our original frozen experiment compares
conventional long short-term memory (LSTM) formulations: S1 uses own-resource
input with independent target models; S2 uses all-resource input with
independent target models; M1 uses all-resource input with a shared
multi-resource model. We evaluate central processing unit (CPU) utilisation, throughput, memory usage
and disk-write activity across TimeTrack development data, Bitbrains
fastStorage and Rnd external environments, and Materna T1--T3 dependent
temporal replication. Post-hoc sensitivity controls separately constrain S2
to M1-selected contexts and increase M1's hidden size from H64 to H129 to
approximately match S2's aggregate parameter count. A matched-context
comparison of Ridge autoregression (AR) and autoregression with cross-resource
inputs (ARX) examines contextual input in a non-LSTM family.
Contextual-input value is selective and model- and environment-dependent;
memory shows repeated benefits, but not universally. Adverse M1-versus-S2
throughput, memory and disk-write point estimates persist in all six blocks
under both separate controls. This supports an accuracy--consolidation
trade-off for the evaluated shared-LSTM design, not a causal attribution to
sharing alone. Ridge confirms that particular contextual effects are
model-specific. Small effects can change sign across seeds, and relative
improvement does not ensure superiority to persistence. Contextual input and
representation sharing therefore require separate, resource-specific
evaluation when consolidating forecasting models.
\end{abstract}

\begin{keywords}
runtime forecasting \sep
cloud telemetry \sep
multi-metric forecasting \sep
multi-task learning \sep
contextual input \sep
shared representations
\end{keywords}

\maketitle

\section{Introduction}
\label{sec:introduction}

Cloud software systems expose multiple streams of runtime telemetry.
Monitoring records past and present resource use; forecasting can supply
advance information for adaptation and resource management
\citep{Taherizadeh2018Monitoring,Bauer2020SelfAware}.
For software engineers, two design decisions follow: should a target's
forecaster include telemetry from other resources, and should several targets
share one predictive representation? The first changes the information
available to a model; the second changes how models are organised and trained.
Consolidating target-specific models can reduce the number of models to
manage, but its predictive consequences need not be favourable for every
resource. We study forecasting accuracy and model consolidation, rather than
the effectiveness of downstream control policies.

Established forecasting research covers several ways of learning from
runtime and resource telemetry. Independent runtime-series benchmarking spans
classical statistical, neural and zero-shot foundation models
\citep{DiMenna2026RuntimeForecasting}. Cloud studies include multivariate
resource inputs, recurrent architectures, joint multi-output prediction and
hybrid models evaluated across multiple provider environments
\citep{Ullah2023CloudBiLSTM,Tripathi2025MMRP}. Global forecasting also learns
parameters across collections of series \citep{Salinas2020DeepAR}.
These designs establish the feasibility of multi-resource prediction, but
do not make additional inputs and model consolidation interchangeable choices.
Our focus is empirical evidence for choosing between these arrangements for
particular resource targets, rather than a new forecasting architecture.

We distinguish \emph{input information sharing} from \emph{predictive
representation sharing}: a target-specific model may use other resources'
histories without sharing learned parameters across targets. A common
representation instead learns several targets jointly. Multi-task learning
offers inductive transfer through related tasks \citep{Caruana1997Multitask},
but task combinations and sharing arrangements require evaluation
\citep{Standley2020TaskGrouping,Fifty2021TaskGroupings,Nochumsohn2025MTLinear}.
Independent all-input target models therefore provide an intermediate
comparison between own-resource forecasting and a shared multi-resource model.

The closest examined designs span shared ten-channel prediction
\citep{Ullah2023CloudBiLSTM}, independently forecast runtime series and
zero-shot foundation models \citep{DiMenna2026RuntimeForecasting}, and MMRP's
univariate/multivariate single-target CPU and memory forecasting
\citep{Tripathi2025MMRP}. MMRP provides relevant prior art for comparing
own-resource and contextual inputs. Among these examined designs, the
remaining gap is a complete within-study comparison that also includes
independent all-input models alongside joint resource forecasting through a
shared representation.

We evaluate conventional LSTM formulations with independent target models
using own-resource histories (S1) or all-resource histories (S2), and a shared
representation with resource-specific heads using all-resource histories
(M1). S2 versus S1 evaluates contextual information, M1 versus S2 the
implemented shared-LSTM transition, and M1 versus S1 the complete joint
system. The evaluated resources are CPU utilisation, throughput, memory usage
and disk-write activity. TimeTrack supplies development evidence, Bitbrains
fastStorage and Rnd provide external environments, and Materna T1--T3 provide
dependent temporal replication. The comparisons retain source-specific
eligible populations and evaluation protocols.
Sections~\ref{sec:framework} and~\ref{sec:methodology} define the original
formulations and evaluation.

The four research questions in Section~\ref{sec:research-questions} concern
contextual input (RQ1), the operational shared-LSTM comparison (RQ2), the
complete joint system (RQ3), and heterogeneity and accuracy--efficiency
trade-offs (RQ4). We interpret them conditionally: which resources and
environments benefit, how sensitive are the contrasts to context selection
and capacity, and which contextual effects recur in another model family?
The original experiment is frozen; subsequent controls test the robustness
of its interpretation rather than replace its primary results.

Two post-hoc LSTM controls examine alternative explanations for the M1--S2
contrast. The common-context control constrains S2 to the frozen contexts
selected by the original M1. The capacity control increases M1's hidden size
from H64 to H129, approximately matching the aggregate parameter count of
S2's independent target models while selecting M1's context and training
epoch afresh using inner validation. These are separate sensitivity controls,
not a jointly context- and capacity-matched factorial intervention.
A further post-hoc comparison uses deterministic Ridge autoregression (AR)
and autoregression with cross-resource inputs (ARX). Both use the same frozen
M1-selected context for each entity and fold, with regularisation selected
only on inner validation. This tests whether contextual-input directions
recur outside the LSTM family under a matched-context design.

The original evidence in Section~\ref{sec:cross-dataset-synthesis} and the
post-hoc controls support selective contextual value. Memory shows repeated
positive patterns, including agreement between LSTM and Ridge in several
external blocks, but particular effects depend on model family and environment.
Throughput, memory and disk-write M1-versus-S2 point estimates remain adverse
in all six blocks under both separate LSTM controls. Context selection and
approximate aggregate capacity matching therefore do not remove those adverse
directions. They do not identify sharing alone as the cause: architecture,
optimisation and training objectives remain part of the implemented contrast.
In the original five-horizon external design, H64 M1 uses approximately
$74.5\%$ fewer parameters and $75\%$ fewer models than S2 for each target set.
This is a model-consolidation trade-off, not evidence of measured production
savings. Seed-sensitive small effects and cases of poor persistence skill
further bound the interpretation of relative gains. We make four contributions:

\begin{enumerate}
  \item We evaluate S1/S2/M1 operational contrasts that distinguish
    contextual-input value from the accuracy consequences of the implemented
    shared-LSTM design and its model consolidation.
  \item We characterise resource-, environment- and model-dependent contextual
    value using the original LSTM comparison and a post-hoc matched-context
    Ridge AR/ARX comparison across all six evidence blocks.
  \item We test the robustness of the shared-LSTM contrast with separate
    common-context and H129 approximate aggregate-capacity controls,
    distinguishing persistent adverse point estimates from causal claims
    about representation sharing alone.
  \item We relate forecasting accuracy to model-management choices while
    bounding conclusions by matched-cell behaviour, existing-seed sensitivity,
    persistence skill and the eligible populations studied.
\end{enumerate}

\section{Related Work}
\label{sec:related-work}

\subsection{Runtime and Cloud-Resource Forecasting}
\label{sec:rw-runtime}

Runtime monitoring supplies observations about software behaviour and the
resources supporting it. In self-adaptive and cloud systems, these observations
inform decisions about configuration, capacity and resource allocation
\citep{Taherizadeh2018Monitoring,Bauer2020SelfAware}. Forecasting extends this
information beyond the current state, potentially allowing adaptation before
demand changes take effect. However, predictive accuracy is only one component
of a management loop: allocation policies, actuation delays and uncertainty
also affect operational outcomes \citep{Qu2018Autoscaling}. Forecast quality
and controller effectiveness therefore require distinct evaluation.

This distinction is reflected in prediction-based provisioning and
software-service research. ARIMA workload prediction has been investigated
together with its consequences for cloud application quality of service,
whereas empirical service-quality forecasting compares the accuracy and cost
of prediction methods \citep{Calheiros2015ARIMA,Syu2017DynamicQoS}. These
studies establish forecasting as a software-engineering concern rather than
solely an algorithmic benchmark. They do not imply that a lower forecasting
error automatically improves service levels. The present study evaluates
telemetry forecasts, not an autoscaling controller, service-level compliance or
deployed resource-allocation policy.

Recent runtime forecasting also extends beyond infrastructure utilisation.
\citet{DiMenna2026RuntimeForecasting} compare classical statistical, neural
and foundation models on 110 daily runtime series from 29 applications,
including crash rates and several application-time measurements. Forecasts
extend to 14 days. The experiments operate independently on each metric:
conventional models are fitted per series, while foundation models are used
zero-shot on individual histories. Multivariate forecasting is left for future
work. This model-family comparison informs runtime prediction, but does not
evaluate sharing among jointly predicted resource tasks.

Cloud-resource forecasting already includes correlated multivariate inputs,
deep recurrent predictors and joint resource outputs
\citep{Tran2018MultivariateFuzzy,Xu2022esDNN,Ullah2023CloudBiLSTM}.
These capabilities establish the resource-forecasting problem addressed here,
rather than an architectural contribution of the present study.

Hybrid architectures, uncertainty-aware predictions and evaluations on multiple
cloud providers further broaden this evidence
\citep{Tripathi2025MMRP,Rossi2025CloudUncertainty,Zhang2025SWTCLSTM}.
Online multivariate resource forecasting and microservice-indicator forecasting
address additional operational settings \citep{Gao2026OMCR,Gao2026MaMiClif}.
Thus multivariate inputs, recurrent resource models, joint forecasting and
multi-provider evaluation are established capabilities, not architectural
contributions of this study.

The TimeTrack dataset places these questions in an OAI CI/CD setting.
\citet{Meliani2025TimeTrack} document its computing and network telemetry and
report contemporaneous cross-resource associations, including CPU--memory
relationships on some machines. Such associations motivate considering other
resource histories, but do not establish incremental predictive value after
conditioning on the target's own history. That distinction requires a
forecasting comparison rather than an inference from correlation alone.

\subsection{Information Scope and Sharing in Forecasting}
\label{sec:rw-information-sharing}

Three design choices must be distinguished: which channels supply historical
inputs, which variables form the prediction targets, and which learned
parameters are common to those targets. A multichannel input can support a
single target; multiple future values can instead represent different horizons
of that target. Neither arrangement by itself establishes resource-task
sharing. Conversely, sharing parameters does not require mixing channels in
each input window.

The closest studies instantiate different parts of this design space
(Table~\ref{tab:related-work-comparison}). In particular, MMRP's
univariate/multivariate single-target comparisons
\citep{Tripathi2025MMRP} are a meaningful contextual-input precedent
related to S1--S2, not evidence of cross-target sharing.

Other designs reinforce the distinction between covariates and targets.
The esDNN study forecasts CPU usage from several resource covariates.
Models that account for uncertainty instead compare predictions for one
resource with joint processing-unit and memory forecasts, alongside transfer
learning
\citep{Xu2022esDNN,Rossi2025CloudUncertainty}. These are different ways of
using several resource series, rather than interchangeable definitions of
multivariate forecasting.

The distinction also matters outside cloud telemetry. PatchTST processes
channels independently while reusing embedding and Transformer weights across
them \citep{Nie2023PatchTST}. Channel independence therefore does not mean
independently fitted models. In contrast, iTransformer represents individual
variates as tokens and models relationships between them through attention
\citep{Liu2024iTransformer}. The relevant difference here is how information
is combined, not a general comparison of Transformer architectures.

Sharing can also occur across series. DeepAR learns recurrent parameters
across collections of related series, while the local/global
forecasting framework distinguishes separately fitted functions from a common
function learned across series
\citep{Salinas2020DeepAR,MonteroManso2021GlobalForecasting}. Such sharing need
not provide every series as simultaneous input channels. Together, these
designs show why input availability, output organisation and shared predictive
representation require separate experimental controls.

\subsection{Multi-Task Learning, Task Relatedness, and Adverse Transfer}
\label{sec:rw-multitask}

Multi-task learning uses relationships among prediction tasks to provide an
inductive bias through shared learning \citep{Caruana1997Multitask}.
Common representations can transfer useful information, but the benefit
depends on task relationships, data and the chosen sharing structure
\citep{Yu2025MultitaskFundamentals}. Learning all available tasks together is
therefore not an unconditional improvement over separate models. Task-grouping
research makes this dependence explicit by selecting combinations under
accuracy and computational constraints, or by estimating how learning one task
affects another \citep{Standley2020TaskGrouping,Fifty2021TaskGroupings}.

Optimisation research addresses a related but distinct issue: joint objectives
can favour different updates. Multi-objective formulations and gradient
modification methods provide alternatives to an unexamined fixed combination
of task losses \citep{Sener2018MultiObjective,Yu2020GradientSurgery}. These
results motivate testing sharing rather than assuming that nominally related
outputs should always use one representation. They do not establish that
gradient conflict, task interference or competition for capacity occurs in a
particular resource forecaster without corresponding measurements.

Forecasting-specific multi-task work already investigates such questions.
NARMAX-LSTM shares LSTM layers across station-level forecasting tasks,
compares single-task baselines and explicitly discusses negative transfer
\citep{Nikentari2022NARMAXLSTM}. MTLinear treats variates or groups as tasks,
analyses gradient relationships and balance, and evaluates grouping and
efficiency \citep{Nochumsohn2025MTLinear}. Cloud workload forecasting has
also been formulated through multi-task learning, with a task construction
different from our resource-head formulation \citep{Bi2019CloudMultitask}.
Negative transfer and selective sharing are thus established concepts. For the
present comparison, degradation after introducing M1 relative to independent
all-input S2 models is operationally consistent with adverse transfer; it does
not by itself identify an optimisation or capacity mechanism.
Neither task grouping nor gradient modification is evaluated in the present
experiment.

\subsection{Evaluation Gap and Study Scope}
\label{sec:rw-gap}

The closest inspected designs cover individual parts of the
information-sharing design space, but not the complete within-study S1/S2/M1
operational comparison (Table~\ref{tab:related-work-comparison}). This
bounded observation is not a literature-wide absence or priority claim.

\input{tables/related_work_comparison}

S2 serves as the intermediate experimental control, rather than introducing
multivariate single-target forecasting. It allows the value of additional
resource histories to be examined separately from the implemented transition
to cross-target representation sharing. S1--S2 and S2--M1 are matched
operational contrasts under common resource, horizon and temporal support,
not pure causal interventions. S1/S2 differ in input width, parameter count
and selected settings; S2/M1 also differ in aggregate capacity, model count,
optimisation and target-specific versus common context selection. The
decomposition therefore characterises the evaluated implementations, not a
parameter-matched or complete factorial effect.

The study retains these contrasts through TimeTrack development evidence,
separate Bitbrains fastStorage and Rnd external environments, and Materna
T1--T3 as dependent temporal replication. Explicit eligible-population
restrictions delimit applicability, while forecasting accuracy and
model-consolidation evidence remain separate outcomes. This scope supports
resource-, environment- and time-dependent assessment without treating
consecutive blocks as independent datasets or extending conclusions to all
shared architectures. Section~\ref{sec:framework} defines the formulations
and the matched comparisons used to address this question.

\section{Multi-Metric Runtime Forecasting Framework}
\label{sec:framework}

The framework distinguishes access to resource histories from sharing a
learned predictive representation across resource targets.
Figure~\ref{fig:sharing-framework} summarises S1, S2 and M1 and their matched
operational contrasts.

\begin{figure*}[t]
  \centering
  \includegraphics[width=\textwidth]{figures/framework_s1_s2_m1.pdf}
  \caption{S1 fits independent target models from own-resource histories;
  S2 supplies all histories to each independently trained target model.
  M1 uses all histories with a shared LSTM representation and resource-specific
  heads. S2 versus S1 assesses contextual input; M1 versus S2 the implemented
  move to representation sharing; M1 versus S1 the complete joint-system
  change. Models are fitted separately for each entity, outer fold and seed;
  S1/S2 use four networks per target set and M1 uses one joint network.
  These are matched operational comparisons, not a complete
  $2\times2$ design, parameter-matched interventions or pure causal
  isolation of either design choice.}
  \label{fig:sharing-framework}
\end{figure*}

\subsection{Problem Formulation}
\label{sec:problem}

Let $\mathcal{R}=\{\mathrm{cpu},\mathrm{thr},\mathrm{mem},\mathrm{disk}\}$
denote the resource set and let $R=|\mathcal{R}|=4$. At discrete time $t$,
the multivariate telemetry observation is
$\mathbf{x}_t=(x_t^{(r)})_{r\in\mathcal{R}}\in\mathbb{R}^{R}$. For a context
length $L\in\mathbb{N}$, the complete history available at forecasting origin
$t$ is
\begin{equation}
  \mathbf{X}_{t,L}
  = [\mathbf{x}_{t-L+1},\ldots,\mathbf{x}_{t}]^{\mathsf T}
  \in \mathbb{R}^{L\times R}.
  \label{eq:multiresource-context}
\end{equation}
Let $\mathcal{H}\subset\mathbb{N}_{>0}$ be a finite set of forecast horizons.
For target resource $r$, $y_{t+h}^{(r)}$ is the future target at horizon
$h\in\mathcal{H}$ and
$\mathbf{y}_{t,\mathcal{H}}^{(r)}=(y_{t+h}^{(r)})_{h\in\mathcal{H}}$ is its
multi-horizon target vector. A forecaster produces
\begin{equation}
  \widehat{\mathbf{y}}_{t,\mathcal{H}}^{(r)}
  = F_r(\mathcal{I}_{t,L};\boldsymbol{\theta}),
  \label{eq:forecasting-task}
\end{equation}
where $\mathcal{I}_{t,L}$ denotes the formulation-dependent information set
and $\boldsymbol{\theta}$ denotes the relevant model parameters.

The notation applies to one entity: a TimeTrack machine or an external VM.
No network pools training observations across entities. For each entity,
outer fold and seed, S1 and S2 fit four independent target-specific networks,
whereas M1 fits one joint four-resource network. Sharing in M1 is across
resource targets within that entity, not across machines or VMs. These
four-versus-one counts describe the local target-set model complement;
candidate-context inner fits and fresh final refits are separate training
runs, not additional architecture families or one globally deployed system.

The context contains only observations available at or before $t$, whereas
each target lies strictly after the forecasting origin. Resource channels are
aligned within a context, and the horizon index is interpreted on the time
grid of the dataset under study. This notation does not require one universal
sampling interval: elapsed-time mappings can differ across datasets provided
that the ordering of observations, origins, and future targets is preserved.
An eligible origin must supply the information required by every formulation
and the complete target vector used in their comparison.

Comparisons among the three formulations use matched forecasting origins,
targets, horizons, and data
partitions. Consequently, the contrasts concern changes in information and
parameter sharing rather than changes in the evaluated sample.

\subsection{Resource Metrics and Multi-Horizon Forecasting}
\label{sec:resource-horizons}

The resource vector comprises CPU utilisation, aggregate throughput, memory
usage, and disk-write activity. CPU utilisation represents processor demand;
throughput represents the aggregate rate of data transferred through the
runtime environment; memory usage represents occupied memory; and disk-write
activity represents write traffic to persistent storage. These channel
definitions are intentionally independent of dataset-specific field names,
units, and sampling intervals. Their concrete mappings are specified with the
experimental data in Section~4.

Every forecasting origin yields one prediction for each
$h\in\mathcal{H}$. The horizon set is therefore a property of an experimental
instantiation rather than a restriction of the framework: $\mathcal{H}$ may be
any finite set of positive offsets supported by the data. The same distinction
applies to $L$. Candidate context lengths and horizon values are fixed before
model evaluation, but their numerical values may vary with sampling cadence
without changing the S1, S2, or M1 definitions. Each formulation predicts its
full horizon vector jointly, rather than fitting a separate model for every
horizon.

The four channels have distinct operational meanings and are not treated as
interchangeable measurements. A resource may nevertheless carry
contemporaneous or lagged information relevant to another target. The
framework makes this a testable modelling question by varying access to the
channel histories while holding the target definition fixed. Joint horizon
prediction similarly holds the output organisation fixed across S1, S2, and
M1, so that information sharing is not confounded with a change from single-
to multiple-horizon training.

Resource channels may differ substantially in scale and distribution. Model
training therefore operates on resource-specific transformed and scaled
targets, while evaluation can return forecasts to their native resource
scales. The transformations and training-partition scaling rules belong to the
experimental instantiation. At the framework level, their purpose is to place
resource errors on comparable optimisation scales, particularly when a single
objective combines several target resources.

\subsection{S1 Local-Resource Forecasting}
\label{sec:s1}

S1 is the local-information formulation. For target resource $r$, its input is
the single-channel sequence
\begin{equation}
  \mathbf{X}_{t,L}^{\mathrm{S1},r}
  = [x_{t-L+1}^{(r)},\ldots,x_t^{(r)}]^{\mathsf T}
  \in\mathbb{R}^{L\times 1},
  \label{eq:s1-input}
\end{equation}
and its output is
\begin{equation}
  \widehat{\mathbf{y}}_{t,\mathcal{H}}^{(r)}
  = f_r^{\mathrm{S1}}
    (\mathbf{X}_{t,L}^{\mathrm{S1},r};\boldsymbol{\theta}_r^{\mathrm{S1}})
  \in\mathbb{R}^{|\mathcal{H}|}.
  \label{eq:s1-forecast}
\end{equation}
A distinct parameter set $\boldsymbol{\theta}_r^{\mathrm{S1}}$ is learned for
each resource; no learned representation is shared across targets. Each S1
model has its own temporal encoder, prediction mapping and optimisation
objective for the selected target resource and horizons. S1 remains
multi-horizon and provides the local-history baseline. Using the same model
family across formulations focuses the comparison on information structure.
S1 withholds non-target channels; it does not assume that resources are
independent.

\subsection{S2 Cross-Resource Contextual Forecasting}
\label{sec:s2}

S2 expands the information set while retaining target-specific modelling. For
every $r\in\mathcal{R}$, the input is the complete history
\begin{equation}
  \mathbf{X}_{t,L}^{\mathrm{S2},r}=\mathbf{X}_{t,L}
  \in\mathbb{R}^{L\times R},
  \qquad
  \widehat{\mathbf{y}}_{t,\mathcal{H}}^{(r)}
  = f_r^{\mathrm{S2}}
    (\mathbf{X}_{t,L};\boldsymbol{\theta}_r^{\mathrm{S2}}).
  \label{eq:s2-forecast}
\end{equation}
Each $\boldsymbol{\theta}_r^{\mathrm{S2}}$ is trained independently.
Non-target histories act as covariates, but the output and loss remain
specific to $r$: gradients from one target cannot update another model's
representation. Thus all-resource input does not establish shared learning
across targets.

Under matched support, S2 versus S1 is a controlled contrast between
all-resource and own-resource histories while preserving the separate-model
organisation. It isolates the empirical contribution of cross-resource
context under the specified design; it does not presume that every additional
channel is informative, nor that the contribution has a particular sign.

\subsection{M1 Shared Multi-Resource Forecasting}
\label{sec:m1}

M1 combines all-resource input with a representation shared across target
tasks. A temporal encoder $E_{\boldsymbol{\theta}}$ maps the complete context
to a latent vector $\mathbf{z}_t$, and one resource-specific head
$g_{r,\boldsymbol{\phi}_r}$ maps that representation to the horizon vector for
resource $r$:
\begin{equation}
  \begin{aligned}
    \mathbf{z}_t
      &= E_{\boldsymbol{\theta}}(\mathbf{X}_{t,L}), \\
    \widehat{\mathbf{y}}_{t,\mathcal{H}}^{(r)}
      &= g_{r,\boldsymbol{\phi}_r}(\mathbf{z}_t),
      && r\in\mathcal{R}.
  \end{aligned}
  \label{eq:m1-forecast}
\end{equation}
The resulting output has resource-by-horizon structure
$\widehat{\mathbf{Y}}_{t}\in\mathbb{R}^{R\times|\mathcal{H}|}$.
The common encoder is updated by all resource losses; head gradients update
only the corresponding head, retaining resource-specific predictive mappings.
This single multi-output model introduces no attention, gating, task routing
or adaptive sharing.

Let $\mathcal{B}$ be a training mini-batch, and let
$\widetilde{y}_{i,h}^{(r)}$ and
$\widehat{\widetilde{y}}_{i,h}^{(r)}$ denote the scaled target and prediction
for sample $i$, resource $r$, and horizon $h$. M1 uses the equally weighted
scaled-space objective
\begin{equation}
  \mathcal{L}_{\mathrm{M1}}
  = \frac{1}{|\mathcal{B}|R|\mathcal{H}|}
    \sum_{i\in\mathcal{B}}
    \sum_{r\in\mathcal{R}}
    \sum_{h\in\mathcal{H}}
    \left(
      \widehat{\widetilde{y}}_{i,h}^{(r)}
      - \widetilde{y}_{i,h}^{(r)}
    \right)^2.
  \label{eq:m1-loss}
\end{equation}
This is equivalently the arithmetic mean of the resource-specific mean squared
errors over the four resources, with horizons weighted equally within each
resource. The targets are transformed and robustly scaled before the losses
are computed. S1 and S2 use the corresponding single-resource scaled mean
squared error. M1 versus S2 compares the implemented joint network with
independent target networks while both receive multichannel input.

Equal arithmetic weights fix the loss coefficients, not the realised loss or
gradient contribution of each resource. Residual distributions and model
derivatives can produce unequal contributions despite these equal nominal
weights. The objective does not learn task weights or conditionally suppress
a resource. Scaling parameters are fitted only within the training partitions
described later.

\subsection{Controlled Information-Sharing Decomposition}
\label{sec:decomposition}

The three matched operational formulations define three comparisons, each evaluated for the same
resource, horizon, and matched support. Let $Q(F)$ denote a lower-is-better
error functional for formulation $F$ under such a fixed evaluation unit. For
a baseline $A$ and candidate $B$, define the signed error reduction
\begin{equation}
  D_Q(A\rightarrow B)=Q(A)-Q(B),
  \label{eq:controlled-reduction}
\end{equation}
so that a positive value indicates lower error for $B$. Exact metric,
aggregation, and percentage-normalisation rules are specified in the
experimental methodology. The framework-level contrasts are
\begin{equation}
  \begin{aligned}
    \mathcal{C}_{\mathrm{context}}
      &= D_Q(\mathrm{S1}\rightarrow\mathrm{S2}), \\
    \mathcal{C}_{\mathrm{sharing}}
      &= D_Q(\mathrm{S2}\rightarrow\mathrm{M1}), \\
    \mathcal{C}_{\mathrm{joint}}
      &= D_Q(\mathrm{S1}\rightarrow\mathrm{M1}).
  \end{aligned}
  \label{eq:controlled-contrasts}
\end{equation}

$\mathcal{C}_{\mathrm{context}}$ compares cross-resource with local-resource
input while retaining independently trained target models.
$\mathcal{C}_{\mathrm{sharing}}$ compares shared with separate predictive
representations while holding all-resource input available to both
formulations. $\mathcal{C}_{\mathrm{joint}}$ compares the complete joint
multi-metric system with the local-information baseline and therefore combines
both changes. The first two contrasts prevent the usefulness of additional
resource streams from being conflated with the usefulness of a shared learned
representation. All three remain empirical matched comparisons: the framework
defines what changes between formulations but imposes no expected ordering or
sign on their observed errors.

Matching requires the same eligible origins, future targets, resource and
horizon coordinates, and evaluation partitions for the formulations entering
a contrast. Results are interpreted resource by resource rather than by
averaging errors expressed in heterogeneous native units. Because the three raw
contrasts are differences of the same error functional $Q$ evaluated on matched
support, the S1-to-M1 raw error difference equals the sum of the S1-to-S2 and
S2-to-M1 raw differences under the sign convention above. This algebraic
identity does not imply that the underlying modelling effects are
mechanistically additive or causally isolated. In particular, the
baseline-normalised percentage gains used for empirical reporting are not
generally additive because each comparison is normalised by its own reference
error. Its role is
to characterise the complete change from local, separate forecasting to the
shared multi-resource system, while the two intermediate contrasts identify
which modelling decision differs.

\section{Experimental Methodology}
\label{sec:methodology}

This section instantiates the controlled formulations from
Section~\ref{sec:framework} across one development dataset, two external
confirmation environments, and three temporal replication blocks.  The design
keeps the information-sharing contrasts fixed while adapting only those
quantities, such as sample offsets and raw field names, that necessarily depend
on the telemetry source.

\subsection{Research Questions}
\label{sec:research-questions}

The evaluation addresses exactly four research questions.  Each question is
operationalised by the matched formulation contrasts defined in
Section~\ref{sec:decomposition}; it is not answered by comparing unrelated
model families or unmatched samples.

\begin{description}
  \item[RQ1:] Does cross-resource context improve forecasting relative to
    local-resource forecasting?  This question is evaluated by comparing S2
    with S1, which changes the input from the target resource alone to all four
    resource histories while retaining separate target-specific models.
  \item[RQ2:] How does the implemented shared-LSTM formulation compare with
    separate contextual target models? This question is evaluated by
    comparing M1 with S2, for which both formulations receive all-resource
    input but only M1 shares a temporal encoder across target tasks.
  \item[RQ3:] Does the complete shared multi-metric system outperform
    local-resource forecasting?  This end-to-end question compares M1 with S1
    and therefore measures the joint change in input scope and representation
    sharing.
  \item[RQ4:] How stable are these effects across resources, environments,
    temporal blocks, horizons, and evaluation units, and what
    accuracy--efficiency trade-offs arise?  Stability is examined without
    collapsing the resource-specific contrasts into heterogeneous native-unit
    errors.
\end{description}

The questions specify estimands rather than expected outcomes.  In particular,
neither access to additional channels nor parameter sharing is assumed to be
beneficial, and efficiency is assessed separately from forecasting accuracy.

\subsection{Datasets}
\label{sec:datasets}

Table~\ref{tab:dataset-summary} summarises the experimental environments and
their roles.  All sources are mapped to the same four resource semantics, but
their sampling grids, eligible populations, and evidential roles remain
distinct.

\subsubsection{TimeTrack Telemetry}
\label{sec:timetrack-data}

TimeTrack is the framework-development source.  It contains telemetry from the
OpenAirInterface (OAI) CI/CD system.  Seven machines were observed between
24 June and 19 July 2024.  The canonical panel
is the exact sorted intersection of timestamps across the six source files; no
nearest-time matching, resampling, interpolation, imputation, or value cleaning
is used.  Its measured median sampling interval is 42.285164 seconds.

Discontinuities are represented explicitly.  A new canonical segment begins
when adjacent retained timestamps are separated by more than 84.57 seconds,
twice the median interval, yielding nine segments.  Historical windows and
future targets cannot cross these boundaries.  The four study targets are
machine-level CPU utilisation, aggregate bidirectional throughput, used memory,
and disk-write activity.  Throughput is constructed for each machine as the sum
of validated \texttt{bond0} transmit and receive rates; used-memory and
disk-write fields are retained at machine level.  TimeTrack supplies internal
development evidence and is not treated as an independent external dataset.

\subsubsection{Bitbrains fastStorage and Rnd}
\label{sec:bitbrains-data}

The Bitbrains fastStorage and Rnd traces are kept as two separate external
confirmation environments.  CPU,
network receive plus transmit throughput, memory usage, and disk-write
throughput are mapped respectively to the four common resource targets.  Rows
that are physically invalid are removed without repair, duplicate timestamps
are rejected, and gaps longer than 450 seconds delimit contiguous segments.

Table~\ref{tab:dataset-summary} distinguishes the structural and scale-eligible
pools. A prospective training-scale amendment required positive
transformed-resource interquartile range in every required training partition.
From each restricted pool, VMs were selected by ascending
SHA-256 rank of a frozen environment-and-identifier key.  The rule is
deterministic and uses neither predictive outcomes nor manual substitution.
The selected cohorts are therefore samples from scale-estimable, structurally
eligible populations; they are not claimed to represent every Bitbrains VM.

\subsubsection{Materna Temporal Traces}
\label{sec:materna-data}

Materna provides a secondary temporal replication.  T1, T2, and T3 are
consecutive blocks from the same source, not independent datasets.  They are
processed independently on their native 300-second grid, with contiguous
segments defined per trace and VM.  The resource mapping is CPU usage
percentage, received plus transmitted network throughput in KB/s, memory usage
in KB, and disk-write throughput in KB/s.

Applying the strict training-scale rule jointly across every trace, fold, training phase,
context, and resource retains 74 VMs, or 18.5930\% of the common structural
pool.  The first 35 VMs under the frozen ascending SHA-256 rank are used in all
three blocks, preserving the same identifiers wherever structurally valid and
permitting matched temporal comparisons.  T1--T3 are analysed first as
separate temporal blocks.  Any combined estimate uses equal block weights;
the blocks are never treated as independent observations in a meta-analysis.

\input{tables/dataset_summary}

\subsection{Target Construction and Preprocessing}
\label{sec:preprocessing}

Dataset-specific fields and units are mapped to the semantic order CPU,
throughput, memory, and disk write.  CPU and memory use the identity transform.
Throughput is the sum of receive and transmit rates, or the frozen equivalent
aggregate in TimeTrack, and uses $\log(1+x)$.  Disk-write activity or throughput
uses the same $\log(1+x)$ transform.  These definitions preserve the meaning of
each resource while accommodating different source schemas and native units.

At each legal forecasting origin, these mappings yield a four-resource history
and future resource-by-horizon target tuple.  Direct multi-horizon prediction is
used: every head estimates all prescribed future offsets in one forward pass,
rather than recursively feeding an earlier prediction into a later horizon.
The target construction is identical across formulations on a given dataset.
Thus S1 receives only the target channel's history, S2 receives all four
histories for the same single target, and M1 receives all four histories and
all four target vectors.  Missing or invalid values remove the affected origin
from common support; they are not repaired differently for one formulation.

For resource $r$, let $g_r$ denote its fixed transform.  Within a legal
training partition, transformed values are robustly scaled as
\begin{equation}
  \begin{aligned}
    \widetilde{y}^{(r)}
      &= \frac{g_r(y^{(r)})-m_r}{q_r}, \\
    m_r &= \operatorname{median}_{\mathrm{train}} g_r(y^{(r)}), \\
    q_r &= \operatorname{IQR}_{\mathrm{train}} g_r(y^{(r)}).
  \end{aligned}
  \label{eq:robust-scaling}
\end{equation}
The interquartile range must satisfy $q_r>0$.  The final protocol uses no
epsilon, fallback scale, or outcome-dependent exception.  Inner scalers are fit
only on inner-training data; after model selection, final scalers are refit on
the complete legal outer-training partition.  Validation and test values never
contribute to scaler parameters.

CPU and memory predictions are returned to native units by inverse robust
scaling.  Throughput and disk-write predictions are inverse-scaled, passed
through $\operatorname{expm1}$, and deterministically projected onto the
non-negative domain.  The same transform, scaler, and inverse mapping are used
for matched S1, S2, and M1 fits.  TimeTrack uses eight output horizons whereas
the external datasets use five; this changes output width but not the
resource-transform semantics.

\subsection{Temporal Support and Leakage Control}
\label{sec:leakage-control}

Forecasting origins are ordered chronologically within the support universe.
An origin is legal only when the maximum candidate context and every required
future target are finite and remain in the same contiguous temporal segment.
Common support is constructed at the largest context, so shorter contexts gain
no additional origins.  S1, S2, and M1 consequently use identical origin IDs,
target IDs, machines or VMs, folds, horizons, and seeds within every
comparison.

Five expanding outer folds are used for each machine, VM, environment, or
trace.  The first half of the ordered legal origins precedes the first test
block; five successive, non-overlapping ten-percent blocks form the outer
tests.  Earlier test periods may enter the training history of a later fold.
For each fold, an otherwise eligible training origin is removed unless its
maximum-horizon target timestamp is strictly earlier than the first outer-test
origin.  This target-time purge prevents a training label from extending into
the test period even when its forecasting origin precedes the boundary.

Within the purged outer-training support, the chronologically latest 15\% of
origins is provisional inner validation.  Earlier inner-training origins are
purged again against the first validation origin using the maximum forecast
horizon.  There is no shuffling.  Scalers, model parameters, context selection,
and epoch selection can use only the relevant legal training and validation
partitions; outer-test observations cannot determine filtering, scaling,
context, epoch count, or fitted parameters.

After context and seed-specific epoch selection, each final model is
reinitialised and refit once on the complete legal outer-training support with
the selected fixed epoch count and a newly fitted outer-training scaler.  It
has no validation or test callback.  Each successful final fit is then
evaluated exactly once on its outer-test block.  These controls separate model
selection from evaluation and ensure that formulation effects are not
attributable to different temporal support.

\subsection{Forecast Horizons and Contexts}
\label{sec:horizons-contexts}

TimeTrack uses sample leads
$\{1,2,4,8,16,32,64,128\}$.  At its measured median interval, these correspond
approximately to 0.705, 1.410, 2.819, 5.638, 11.276, 22.552, 45.104, and
90.208 minutes.  Candidate context lengths are 32, 64, and 128 samples,
covering approximately 22.6, 45.1, and 90.2 minutes.

Bitbrains and Materna use five-minute sampling, so their sample leads are
$\{1,2,4,9,18\}$, corresponding exactly to 5, 10, 20, 45, and 90 minutes.
Their candidate contexts are 4, 9, and 18 samples, or approximately 20, 45,
and 90 minutes.  The integer offsets therefore differ by dataset.  The
adaptation preserves a comparable real-time forecast and context range rather
than pretending that one sample index has the same elapsed-time meaning on all
grids.

\subsection{Model Architecture and Training}
\label{sec:architecture-training}

In the original frozen experiment, all three formulations use the same plain recurrent backbone: a two-layer
unidirectional LSTM with hidden size 64, dropout 0.1, and the final recurrent
output as its temporal representation.  S1 has one input channel and one
resource-specific horizon head; S2 has four input channels but remains one
independent model per target; M1 has four inputs, one shared encoder, and four
independent resource heads.  S1 and S2 minimise target-specific scaled mean
squared error, whereas M1 minimises the arithmetic mean of the four
resource-specific scaled mean squared errors. Inner early stopping minimises
validation MSE in this transformed/robust-scaled space and records a best
epoch separately for each seed.

Training uses Adam with learning rate 0.001, batch size 256, at most 50 epochs,
early-stopping patience 7 during inner fitting, gradient-norm clipping at 1.0,
and seeds 0, 1, and 2.  No architecture or optimiser search is performed.  The
plain backbone is deliberately simple because the experiment compares the
operational formulations rather than proposing an architecture-specific
innovation.

Output width affects parameter counts.  For the five-horizon external design,
S1 contains 50,757 parameters per resource model (203,028 across four), S2
contains 51,525 per target model (206,100 across four), and M1 contains 52,500
in its single shared model.  For eight-horizon TimeTrack, the corresponding
counts are 50,952 per S1 model (203,808 across four), 51,720 per S2 model
(206,880 across four), and 53,280 for M1. These H64 counts refer to one
entity/fold/seed target complement, not parameters pooled across the corpus.
The four-versus-one reduction in model count and the aggregate parameter
reduction are architecture inventories; they are not total training-fit
counts. Each candidate context and seed requires its own inner fit, followed
by a fresh final refit for each selected model and seed. Accuracy and model
footprint are therefore reported separately rather than assuming equal
parameter budgets.

Reported fitting time is the sum of successful unique worker fit-hours across
inner candidate fits and final refits. Inner durations include training-loop
and validation-forward time; final durations cover the training loop.
Unsuccessful attempts, skipped/resumed encounters, model and window
construction, preprocessing, checkpoint serialisation, final prediction and
inversion, and orchestration are excluded. Summed worker fit-hours are not
elapsed wall time, energy, deployment cost or production latency. Differences
between S1 and S2 can reflect executed epochs and selected context lengths;
input dimensionality alone does not explain their timing order.

\subsection{Context and Model Selection}
\label{sec:model-selection}

In the original LSTM experiment, context length is selected from the frozen
candidate grid while architecture and optimiser settings remain fixed.
For S1 and S2, one context
is selected independently for each environment or trace, machine or VM, outer
fold, and target resource. The selection score is validation MAE of the
transformed/robust-scaled targets and predictions, evaluated at each seed's
MSE-selected best epoch, then averaged equally over horizons and the three
seeds. M1 selects one context per environment or trace, machine or VM, and
fold by taking the equal mean of the same scaled-space validation MAE over
the four resources, all horizons and all seeds. Exact ties are resolved
in favour of the shorter context.

Early stopping and context selection consequently use different error
functionals in the same training space: scaled MSE chooses the epoch, and
scaled MAE scores the context at that epoch. This context score is not native
evaluation NMAE. For log-transformed targets, inverse transformation and
non-negativity projection need not preserve candidate rankings.
Once a context has been selected, the final fit uses that context, the
same seed, and its own fixed selected epoch.  It starts from a fresh
initialisation, refits the resource scalers on all legal outer-training data,
and performs no test-based early stopping.  Thus context selection does not
stand in for broad architecture or hyperparameter tuning, and the final model
does not inherit inner-fit weights.

The selection units differ deliberately between independent and shared
formulations.  S1 and S2 may choose target-specific context lengths because
each target has its own model.  M1 must choose one common context because a
single shared encoder serves all four targets.  In both cases, the selection
criterion is computed before outer-test evaluation and retains equal horizon
and seed contributions.

\subsection{Evaluation Metrics and Statistical Analysis}
\label{sec:evaluation}

Let $y_i$, $\widehat{y}_i$, and $\overline{y}$ denote the observed value,
forecast, and observed mean for $n$ matched cases of one resource and horizon.
Evaluation uses
\begin{equation}
  \begin{aligned}
    \operatorname{MAE} &= \frac{1}{n}\sum_{i=1}^{n}|y_i-\widehat{y}_i|, \\
    \operatorname{RMSE} &=
      \sqrt{\frac{1}{n}\sum_{i=1}^{n}(y_i-\widehat{y}_i)^2}, \\
    R^2 &= 1-
      \frac{\sum_{i=1}^{n}(y_i-\widehat{y}_i)^2}
           {\sum_{i=1}^{n}(y_i-\overline{y})^2}, \\
    \operatorname{NMAE}_r &= \frac{\operatorname{MAE}_r}{q_r^{\mathrm{native}}},
  \end{aligned}
  \label{eq:evaluation-metrics}
\end{equation}
where $q_r^{\mathrm{native}}$ is the native-unit outer-training IQR for resource
$r$, distinct from the transformed training IQR $q_r$ in
Equation~\ref{eq:robust-scaling}. Primary evaluation follows inverse
transformation and the prescribed non-negativity projection. MAE and RMSE are
reported in each resource's native units; $R^2$ and NMAE are unitless.  Native
errors are never averaged across different resources.

Persistence supplies a separate absolute forecasting sanity check:
\begin{equation}
  \widehat{y}_{t+h}^{(r),\mathrm{pers}}=y_t^{(r)},
  \qquad
  \operatorname{Skill}_{r,h}
  =1-\frac{\operatorname{MAE}_{r,h}^{\mathrm{model}}}
           {\operatorname{MAE}_{r,h}^{\mathrm{pers}}}.
  \label{eq:persistence-skill}
\end{equation}
Persistence skill does not replace the controlled formulation comparisons:
relative improvement over another learned formulation does not imply absolute
superiority over persistence. Reported overall skill averages the existing
per-cell relative skills over entity--fold cells and then horizons; it is not
$1-\operatorname{MacroMAE}_{\mathrm{model}}/
\operatorname{MacroMAE}_{\mathrm{pers}}$. Undefined zero-denominator skills
are excluded from this ancillary mean with their counts retained, not
replaced by zero.

The primary aggregation begins by averaging stochastic-model MAE across the
three seeds within each matched environment or trace, entity, fold, resource,
horizon, and formulation cell.  For a fixed resource and horizon, those
seed-averaged entity--fold cells receive equal weight to form MacroMAE. A
resource-overall MacroMAE then gives equal weight to every horizon.  For
candidate $C$ and baseline $B$, the reported gain is
\begin{equation}
  G(C\!:\!B)=100\,
  \frac{\operatorname{MacroMAE}(B)-\operatorname{MacroMAE}(C)}
       {\operatorname{MacroMAE}(B)}.
  \label{eq:macro-gain}
\end{equation}
A positive value indicates lower candidate error.  The three contrasts are S2
versus S1, M1 versus S2, and M1 versus S1.  Secondary diagnostics include
pooled prediction-level gain, mean and median matched-cell relative gain, and
the fractions or counts of positive VMs, folds, horizons, and seed aggregates.

Prospectively frozen directional classifications combine overall gain
magnitude with horizon, entity, fold and seed consistency. A positive
aggregate has gain strictly greater than zero; exact zero is not positive.
Entity signs aggregate over folds and horizons, fold signs over entities and
horizons, and seed signs over entities, folds and horizons. The following
conditions are conjunctive, with overall gain expressed in percent:
\begin{description}
  \item[\textsc{Broad Positive}.] Overall gain is at least $+1\%$, with at
    least $6/8$ horizons, $5/7$ machines, $4/5$ folds and $2/3$ seeds positive
    in TimeTrack. Externally, the horizon and entity thresholds are $4/5$
    and $25/35$ VMs; fold and seed thresholds are unchanged.
  \item[\textsc{Conditionally Positive}.] The broad-positive rule fails,
    overall gain is strictly positive, and at least $4/8$ horizons, $3/7$
    machines, $3/5$ folds and $2/3$ seeds are positive in TimeTrack.
    Externally, the horizon and entity thresholds are $3/5$ and $15/35$ VMs.
  \item[\textsc{Broad Negative}.] Overall gain is at most $-1\%$, with at
    most $2/8$ horizons, $2/7$ machines, $1/5$ folds and $1/3$ seeds positive
    in TimeTrack. Externally, the horizon and entity limits are $1/5$
    and $10/35$ VMs.
  \item[\textsc{Neutral or Mixed}.] Every remaining case, including point
    gains whose magnitude or lower-level consistency fails these rules.
\end{description}
These classes are not labels assigned solely by the sign of the primary
point estimate. Bootstrap intervals quantify a separate form of uncertainty:
neither interval exclusion of zero nor a significance test enters the
classification predicates. Full trace-level and RQ1--RQ4 predicates,
including precedence and temporal classification-change rules, remain in the
supplementary decision criteria and frozen protocol. No threshold is adapted
after observing outcomes.
For external uncertainty, a paired VM-cluster bootstrap samples 35 VMs with
replacement, retains all five folds and the paired formulation structure, and
operates on seed-averaged matched MAEs.  It uses 10,000 replicates, percentile
95\% intervals (2.5th and 97.5th percentiles), and random seed 20261007.  The
bootstrap is performed independently within each Bitbrains environment and
each Materna trace.  Materna blocks are not resampled as independent datasets,
and this exact external bootstrap is not retrospectively imposed on earlier
exploratory TimeTrack analyses.

\subsection{External and Temporal Replication Design}
\label{sec:replication-design}

The study assigns distinct prospective roles to each source.  TimeTrack is the
framework and development study.  Bitbrains fastStorage and Rnd are separate
external confirmation environments; each is analysed before any
equal-environment synthesis.  Materna is a secondary temporal replication in
which T1, T2, and T3 probe stability across consecutive blocks from one source.
This progression holds the three S1/S2/M1 operational formulations, semantic resource mapping,
controlled contrasts, matched-support principle, and approximately comparable
real-time forecast range constant while allowing source-specific field names,
units, and sample offsets.

No numeric score pools all dataset roles.  Estimates for each Bitbrains
environment precede an unweighted equal-environment summary where applicable.
Materna likewise reports each trace first.  Its descriptive synthesis for any
fixed comparison and resource is
\begin{equation}
  G_{\mathrm{temporal}}
  = \frac{G_{\mathrm{T1}}+G_{\mathrm{T2}}+G_{\mathrm{T3}}}{3}.
  \label{eq:temporal-synthesis}
\end{equation}
This equal-temporal-block estimate is not weighted by rows, origins,
predictions, duration, or fit count.  It is explicitly not an
independent-dataset meta-analysis.  Cross-source synthesis is therefore
descriptive and based on direction and frozen classification consistency,
rather than a pooled error or nominal interval that would erase the sources'
different roles and dependence structures.

\subsection{Population Restrictions}
\label{sec:population-restrictions}

External eligibility requires both structural validity and non-degenerate
training variability.  Specifically, every transformed CPU, throughput,
memory, and disk-write series must have IQR strictly greater than zero in every
required scaler-fit partition.  The partitions cover each fold, inner and final
training phase, and candidate context; for Materna the same requirement also
holds in every temporal block.  There is no epsilon, fallback, resource mask,
or formulation-specific exemption.

Table~\ref{tab:dataset-summary} reports the structural, scale-eligible and
executed populations. Deterministic hash ranking follows eligibility
filtering; selection is neither random nor performance-ranked. Accordingly, external
conclusions concern eligible dynamic, scale-estimable, multi-resource VM
populations.  They do not automatically generalise to excluded idle,
quiescent, resource-constant, intermittently observed, or otherwise
structurally invalid VMs, nor to every VM in the source archives.  This
qualification is part of the estimand itself: deterministic selection makes the
executed cohort reproducible, but it does not convert the cohort into a
probability sample of the rejected or unobserved populations.

\subsection{Post-Hoc Sensitivity Analyses}
\label{sec:posthoc-controls}

Three completed analyses address reviewer challenges after the original
experiment was frozen. They are explicitly \emph{post-hoc sensitivity
controls}, not amendments to the original prospective design or replacements
for its primary results. All retain the frozen cohorts, outer folds,
forecast horizons, purging, legal support, source-specific transforms,
training-only scalers and native-unit evaluation. Models remain
entity-specific, with no cross-machine or cross-VM pooling. Results are
reported separately by block and resource.

\paragraph{S2 with common context.}
For each entity and outer fold, the frozen original H64-M1 context is imposed
on every independent S2 target. Each seed retains the S2 MSE-selected epoch
from its existing inner fit at that imposed context. Exact matching final
artifacts are reused only after identity, support, scaler and hash checks;
otherwise the selected fixed-epoch outer-training refit is performed.
This removes S2's target-specific context-selection freedom relative to
H64 M1, but does not equalise parameter count, architecture or optimisation.

\paragraph{M1 with approximately matched aggregate capacity.}
The shared LSTM hidden size increases from H64 to H129, giving 206,420
parameters externally and 207,980 for TimeTrack, approximately matching the
corresponding four-network S2 complements of 206,100 and 206,880. All other
scientific settings are retained. M1's original inner context and
seed-specific epoch selection are repeated afresh, followed by fresh final
refits; H64 choices and fitted weights are not reused. This controls the
approximate aggregate parameter budget, not task-specific capacity,
architecture or optimisation equivalence. Common-context S2 remains aligned
to H64-M1 contexts, not necessarily the newly selected H129 contexts; the
two LSTM controls are separate, not a jointly matched intervention.

\paragraph{Matched-context Ridge AR/ARX.}
AR uses only the target resource's lag window; ARX uses all four resource lag
windows. Both use the same frozen H64-M1 context per entity and fold, without
context tuning, and directly predict the complete horizon vector for one
target. The existing multi-output Ridge implementation flattens each window
and applies an additional lag-feature StandardScaler fitted only on training
data. Targets retain the source-specific transformed/robust-scaled space.
For each formulation and target, $\alpha\in\{0.01,0.1,1,10,100\}$ is selected
strictly within the inner split using native-unit validation MAE after inverse
transformation/projection, averaged equally over horizons; exact ties favour
the smaller alpha. The selected model is refitted on outer-training data and
evaluated once on outer test. Ridge is deterministic, so no artificial seed
repetitions are introduced. This tests input-scope sensitivity within a
second model family, not equality of LSTM/Ridge selection objectives or a
causal isolation of model family from the specified context and model design.
These post-hoc analyses introduce no additional significance test or pooled
cross-block performance statistic.

\section{Results}

Results are reported as the primary overall MacroMAE gain defined in
Section~4.8. Positive values mean that the candidate formulation has lower
MacroMAE than its baseline; negative values therefore indicate degradation.
Figure~\ref{fig:primary-gains} gives the complete resource-level primary result
for every development, external, and temporal evidence block. The prospective
VM-cluster intervals are used selectively below where they distinguish a
directionally clear effect from an uncertain near-zero result. TimeTrack did
not use the same prospective external bootstrap and is therefore reported
without those intervals.

The three evidence roles are kept distinct throughout this section. TimeTrack
provides the development result, the two Bitbrains environments provide the
external-environment results, and the three Materna traces provide the
prospective temporal results. Each reported gain retains the aggregation and
baseline fixed for its own evidence block. Comparisons between blocks are
therefore descriptive: no prediction-level pooling is used, and the six blocks
are not treated as statistically independent datasets. The frozen qualitative
classes summarise consistency across the prescribed lower-level units, while
the VM-cluster intervals quantify uncertainty for the external and temporal
blocks. Neither replaces the primary resource-overall MacroMAE gain.

The original frozen S1/S2/M1 estimates remain the primary results.
The completed controls and existing-artifact diagnostics below are explicitly
\emph{POST-HOC}; they do not replace those estimates or their frozen classes
and intervals. Table~\ref{tab:posthoc-controls} places the post-hoc results
beside unchanged frozen reference gains, using the designs in
Section~\ref{sec:posthoc-controls}.

\begin{figure*}[t]
\centering
\includegraphics[width=\textwidth]{figures/results_signed_matrix.pdf}
\caption{Primary overall MacroMAE percentage gains for (a) S2 versus S1,
(b) M1 versus S2 and (c) M1 versus S1 across the evaluated evidence
blocks. Positive values indicate lower MacroMAE for the candidate
formulation; negative values indicate higher MacroMAE. Continuous signed
colour uses a symmetric $\operatorname{asinh}(G)$ mapping, with $G$ in
percentage points, as a nonlinear visual aid only. No values are clipped;
the printed raw percentages at the reported precision are authoritative.
Legend ticks are raw percentages. Colour implies neither significance nor
a consistency class. Materna T1, T2 and T3 are dependent temporal blocks.
Comparisons are bounded to the evaluated implementation and eligible
populations.}
\label{fig:primary-gains}
\end{figure*}

\subsection{RQ1 Cross-Resource Context}
\label{sec:results-rq1}

RQ1 compares independent all-resource S2 models with own-resource S1
models. Figure~\ref{fig:primary-gains}(a) shows selective
contextual value, not uniform improvement. In TimeTrack, memory was the sole
positive resource-level effect, while all four frozen classifications remained
neutral or mixed because the required consistency thresholds were not met.

The Bitbrains environments did not reproduce one common direction. In
fastStorage, memory was broadly positive and disk write conditionally positive;
CPU and throughput were neutral or mixed. The memory bootstrap interval
excluded zero, whereas the near-zero throughput interval crossed it. In Rnd,
CPU and memory were conditionally positive, while throughput and disk write
were neutral or mixed. All four Rnd intervals crossed zero, including the
interval for its negative throughput point estimate.

The Materna blocks further changed the resource pattern. T1 memory was
approximately neutral and slightly negative ($-0.02\%$); only throughput
reached the conditionally positive class. The T1 CPU interval was entirely
negative, while the other three intervals crossed zero. T2 CPU and memory
were broadly positive with intervals above zero; throughput and disk write
remained neutral or mixed. T3 memory and disk write were conditionally
positive, but all four T3 intervals crossed zero.

Memory therefore has the clearest repeated positive contextual pattern, but
not universally. CPU, throughput and disk write change direction across
blocks. The resource ordering also changes: the strongest contextual benefit
does not remain attached to one resource, while Rnd throughput and T1 disk
write illustrate adverse effects on resources that are near zero or positive
elsewhere. Contextual value varies by resource, environment and time, not
merely in magnitude around a stable ordering.

\paragraph{Post-hoc cross-model and seed checks.}
The matched-context Ridge ARX:AR comparison
(Table~\ref{tab:posthoc-controls}B) also shows selective contextual value,
but individual directions depend on model family and environment. Memory
improves under Ridge in all five external blocks; fastStorage, Rnd and
Materna T2/T3 have positive point gains in both model families. TimeTrack
memory instead reverses under Ridge, and Materna T1's near-zero negative
LSTM estimate contrasts with a positive Ridge gain. External memory is thus
the strongest recurring favourable contextual pattern, not a universal
benefit. Existing-seed diagnostics find sign changes in 9/24 LSTM S2:S1
block/resource summaries across the three fixed seeds. Small point gains or
sign reversals should not be treated as stable improvements or temporal
regime changes.

The Ridge TimeTrack throughput gain of $-838.45\%$ is diagnosed as
\textsc{Concentrated Degradation}, not broad Ridge failure. The saved-artifact
audit attributes $99.93\%$ of positive excess error to machine06/fold3.
ARX nevertheless improves 21/35 machine/fold pairs, with median ARX/AR MAE
ratio approximately $0.9963$. Predictions and primary MAEs are finite, and
transforms, scalers and support match the frozen design; no implementation
defect is demonstrated. The aggregate penalty must be read with this
concentration diagnosis rather than as a typical-machine effect.

\subsection{RQ2 Operational Shared-Representation Comparison}
\label{sec:results-rq2}

RQ2 compares M1 with S2, holding the all-resource input scope fixed and changing
whether the predictive representation is shared. Throughput, memory and disk
write had negative primary gains in all six blocks
(Figure~\ref{fig:primary-gains}(b)). CPU was positive only in TimeTrack,
fastStorage and Materna T1. These small or near-zero exceptions do not support
uniformly beneficial sharing within the evaluated M1 design.

In TimeTrack, memory was broadly negative; the remaining frozen resource
classifications were neutral or mixed. In fastStorage, CPU was neutral or
mixed and its bootstrap interval crossed zero. Throughput, memory and disk
write were broadly negative with intervals below zero. Rnd was broadly
negative for every resource, and each interval also lay below zero.

Materna T1 CPU was neutral or mixed with an interval crossing zero, whereas
the other resources were broadly negative with intervals wholly below zero.
T2 was broadly negative for all four resources, with entirely negative
overall intervals. In T3, throughput, memory and disk write were broadly
negative; CPU was neutral or mixed under the frozen consistency rule despite
its negative primary gain. All four T3 intervals excluded zero, with the
disk-write interval spanning $[-239.63\%,-5.94\%]$.

Within the controlled shared-LSTM formulation, degradation therefore recurred
outside the CPU exceptions, but its magnitude remained resource- and
block-dependent. External and temporal effects for the other resources often
combined sizeable degradation with a broadly negative class and an interval
below zero. A negative point estimate and a consistently negative result
remain distinct under the frozen protocol.

\paragraph{Post-hoc sharing controls and matched cells.}
Table~\ref{tab:posthoc-controls}A shows H64 M1 against S2 constrained to the
frozen H64-M1 contexts, and H129 M1 against original S2. Throughput, memory
and disk-write point estimates remain adverse in all six blocks under both
separate controls. Context-selection freedom and the original aggregate
capacity difference are therefore incomplete explanations within this
pipeline. The controls do not identify sharing alone as the cause:
architecture, optimisation and training objectives remain part of the
implemented comparison, and common-S2 contexts align H64, not necessarily
H129. Existing-seed diagnostics find sign changes in 4/24 original M1:S2
block/resource summaries, concentrated in the small CPU exceptions and
TimeTrack throughput. Adverse aggregate directions do not imply adverse
outcomes for every seed.

The post-hoc matched-cell diagnostics further separate direction from
magnitude. A cell is entity $\times$ outer fold $\times$ forecast lead, with
MAE first averaged across existing seeds. In all 24 original M1:S2
block/resource summaries, the median candidate/reference MAE ratio exceeds
one and the candidate win rate is below $50\%$. This includes the small
positive CPU MacroMAE exceptions. Materna T1/T3 disk-write median ratios are
approximately $1.0246/1.0347$, despite MacroMAE gains of
$-27.46\%/-71.30\%$. The adverse typical-cell direction is widespread, but
the much larger mean-error penalties do not describe every VM or the
magnitude of degradation in a typical cell.

\input{tables/posthoc_controls}

\subsection{RQ3 Full Joint-System Performance}
\label{sec:results-rq3}

RQ3 compares the complete M1 system with local-resource S1 models.
Figure~\ref{fig:primary-gains}(c) is negative in every displayed cell except
Materna T2 CPU, so the full joint formulation generally did not outperform S1.
In TimeTrack, memory and disk write were broadly negative; CPU and throughput
were neutral or mixed under the prospective consistency criteria.

In fastStorage, CPU and throughput were broadly negative, whereas memory and
disk write were neutral or mixed. The throughput interval was entirely below
zero, while the other three overall intervals crossed zero. In Rnd, every
resource was broadly negative and every interval lay below zero.

In Materna T1, throughput and disk write were broadly negative; CPU and memory
were neutral or mixed, although all four VM-cluster intervals lay below zero.
The T2 CPU exception ($+1.75\%$) was conditionally positive, but its interval
crossed zero ($[-3.52\%,+6.07\%]$). T2 disk write was broadly negative, and
throughput and memory were neutral or mixed. In T3, CPU, throughput and disk
write were broadly negative, while memory was neutral or mixed; only the
memory interval narrowly crossed zero. These differences in classification
and uncertainty remain important despite the nearly uniform negative sign
pattern.

The complete-system comparison incorporates both the move from local to
all-resource inputs and the move from separate to shared representation. Its
observed gain is therefore reported directly rather than reconstructed by
adding the RQ1 and RQ2 percentages. The block-level evidence is not pooled or
treated as six statistically independent datasets.

\subsection{RQ4 Resource, Environment, and Temporal Heterogeneity}
\label{sec:results-rq4}

Figure~\ref{fig:primary-gains} reveals heterogeneity at each prespecified
level. Resource heterogeneity is evident in the selective S2-versus-S1
benefits and changing resource ordering. CPU, throughput and disk write change
direction across evidence blocks, whereas memory is positive in five and
nearly zero and negative only in Materna T1. Repeated positive memory effects
therefore do not imply invariant resource behaviour.

External-environment heterogeneity is stronger for representation sharing: fastStorage
CPU was near zero while the other three resources were negative; Rnd was
broadly negative for all four. Temporal heterogeneity within the dependent
Materna blocks is distinct: M1-versus-S2 CPU changed from positive in T1 to
negative in T2--T3, moving from neutral or mixed to broadly negative and back
to neutral or mixed. Throughput, memory and disk write were broadly negative
in all three Materna blocks, with varying magnitudes. For M1 versus S1, CPU
alone reversed sign in T2; throughput and disk write stayed negative and
memory stayed negative but neutral or mixed across T1--T3.

Sign, magnitude, and frozen class capture different aspects of this
heterogeneity. For example, T3 S2-versus-S1 throughput was slightly positive at
$+0.04\%$ but remained neutral or mixed, while T2 CPU was both positive and
broadly positive. T3 M1-versus-S2 CPU had a negative point estimate and an
interval below zero, yet remained neutral or mixed under the separate
consistency rule. These combinations prevent near-zero sign changes from being
treated as equivalent to changes supported across the prescribed units. They
also show why the within-Materna decision uses the frozen classification and
trace rules alongside the raw gains.

The frozen within-Materna RQ4 outcome was recorded as follows.

\noindent Temporal heterogeneity confirmed. Mechanically, seven of the 12
resource--comparison pairs changed frozen classification across T1--T3 and
five changed raw sign. The S2-versus-S1 trace decision also changed: T1 met the
frozen non-support rule, whereas T2 and T3 were mixed. The M1-versus-S2
trace-level decision was consistently not supported, even though its CPU
classification varied. These prespecified changes satisfy the RQ4 criterion;
no post-hoc temporal regimes or causal mechanisms are inferred.

\subsection{Persistence Boundary}
\label{sec:results-persistence}

Relative formulation gains do not establish superiority to a persistence
forecast. In TimeTrack, S1, S2, and M1 all had positive mean persistence skill
for throughput and disk write, and their CPU means were also positive. Memory
was the boundary: mean skill was $-0.1272$ for S1, $-0.1494$ for S2, and
$-0.6047$ for M1, with positive skill at only two, four, and zero of eight
leads. Thus the large negative M1-versus-S2 memory result coincided with weak
absolute performance relative to persistence, whereas the throughput and
disk-write formulation comparisons occurred in resources with positive
average absolute skill.

Neither Bitbrains environment contained a formulation that beat persistence
for every resource. Overall CPU skill was negative for S1, S2, and M1 in both
environments. In fastStorage, throughput skill was also negative for every
formulation; memory and disk write were positive for all three. Rnd again had
positive memory and disk-write skill for all formulations, while throughput
was only slightly positive for S1 ($0.0384$) and negative for S2 ($-0.0257$)
and M1 ($-0.1176$). Consequently, a positive S2-versus-S1 gain for Rnd CPU or a
relative degradation under M1 does not by itself describe absolute forecasting
skill.

Materna showed a different aggregate boundary. S2 beat persistence in all 12
trace--resource overall aggregates. No formulation, however, beat persistence
in every individual trace--resource--lead cell. At the overall level, M1 failed
the check for T1 throughput ($-0.0608$), T1 disk write ($-1.7590$), T3
throughput ($-0.2629$), and T3 disk write ($-1.3493$); it remained positive for
all four T2 resources. S1 additionally failed only T1 throughput at the overall
level, while S2 had no overall failures. These persistence results delimit
absolute forecasting performance and remain separate from the controlled
S1--S2--M1 comparisons.

Three boundaries follow directly from these observations. First, a positive
gain answers only whether one learned formulation improves on another; it does
not answer whether either improves on persistence. Second, an overall positive
skill does not require every horizon-level cell to be positive, as the Materna
lead-level result demonstrates. Third, absolute skill remains resource
dependent: memory and disk write were the stronger Bitbrains cases, whereas
CPU remained below persistence despite some positive relative formulation
gains. Reporting the persistence check beside the formulation comparison
therefore prevents a relative improvement in an absolutely weak setting from
being described as general forecasting superiority.

\subsection{Accuracy-Efficiency Trade-Off}
\label{sec:results-efficiency}

The original H64 M1 consolidates the target-model complement for one entity,
outer fold and seed. In the five-horizon external design, S1 uses four
networks and 203,028 parameters, S2 four networks and 206,100 parameters, and
M1 one joint network and 52,500 parameters: approximately $74.5\%$ fewer
parameters and $75\%$ fewer models than S2. TimeTrack's eight-horizon
complements contain 206,880 S2 parameters and 53,280 M1 parameters. These
are local architecture inventories, not a globally deployed four-versus-one
system or total training-fit counts. The post-hoc H129 comparison changes
the parameter budget and must not inherit H64's parameter-reduction claim.

\input{tables/efficiency_summary}

The frozen Materna totals are 85.594240 worker fit-hours for S1, 75.303186 for
S2 and 20.405083 for M1. As defined in Section~\ref{sec:architecture-training},
these sum successful inner-candidate and final-refit durations, including
inner validation, across the three traces. They are not elapsed wall time,
energy, deployment cost or production latency. Different executed epochs and
selected contexts can contribute to the S1--S2 order; input width alone does
not explain it. The measured training workload and local model consolidation
coexist with the accuracy penalties above, without establishing production
savings or an architecture-independent efficiency ranking.

\section{Cross-Dataset Synthesis}
\label{sec:cross-dataset-synthesis}

Table~\ref{tab:cross-study-synthesis} retains the frozen source summaries:
TimeTrack development evidence, equal-environment Bitbrains summaries and
equal-temporal-block Materna summaries. These are not combined into a
six-block global score. The \emph{post-hoc} controls and diagnostics in
Table~\ref{tab:posthoc-controls} and Section~5 qualify their interpretation
without replacing the primary gains, frozen classes or intervals. Agreement
concerns direction under the evaluated designs, not independent replication
of a uniquely identified sharing mechanism.

\input{tables/cross_study_synthesis}

\subsection{TimeTrack}

The development result separates weak, selective contextual effects from
the generally adverse operational shared-LSTM contrast. The small positive
S2:S1 memory and M1:S2 CPU point estimates are seed-sensitive. Ridge reverses
the memory-context direction, while its extreme throughput penalty is
concentrated in one machine/fold rather than typical of all machines
(Section~\ref{sec:results-rq1}). Both separate LSTM controls retain adverse
throughput, memory and disk-write point estimates. TimeTrack establishes
conditional design boundaries, not independent external confirmation.

\subsection{Bitbrains External Confirmation}

fastStorage and Rnd provide distinct external environments. Memory has
favourable contextual point agreement between LSTM and Ridge in both, whereas
throughput reverses direction between the model families in each environment.
The common-context and H129 controls both preserve adverse M1:S2 throughput,
memory and disk-write point estimates. These block-specific results separate
the recurring external-memory signal from environment- and model-dependent
contextual effects. The frozen equal-environment summary is retained without
using it to erase those differences.

\subsection{Materna Temporal Replication}

T1--T3 are consecutive dependent blocks using the same selected VM
identifiers, not three independent datasets. LSTM and Ridge agree favourably
on memory in T2/T3; T1's near-zero, seed-variable LSTM point estimate differs
from Ridge's positive direction. CPU, throughput and disk-write context
directions remain conditional on block and model family. For the operational
sharing comparison, both controls retain adverse throughput, memory and
disk-write point estimates in every block. T1/T3 disk-write magnitudes require
the matched-cell and concentration qualifications in
Section~\ref{sec:results-rq2}, rather than a claim of uniformly severe VM
degradation.

\subsection{Generalisation Boundaries}

\paragraph{Contextual inputs.}
External memory is the strongest recurring favourable case, but particular
contextual directions do not transfer uniformly between LSTM and Ridge.
Existing-seed sign variation also limits claims based on small primary gains.
The evidence supports validation by resource, environment and model family,
not universal benefit or universal harm from additional telemetry.

\paragraph{Shared representation.}
The adverse throughput, memory and disk-write point directions survive both
separate controls, so neither the original capacity difference nor context
selection alone removes the observed pattern. The original matched-cell
medians and win rates reinforce an adverse typical-cell direction even where
a small CPU MacroMAE gain is positive. Architecture and optimisation remain
part of the contrast; sharing alone is not causally identified.

\paragraph{Complete joint system and consolidation.}
The frozen M1:S1 gains are negative in every displayed resource/block
combination except Materna T2 CPU. This repeated direction is reported without
a six-block arithmetic mean or another pooled global statistic. The local
model complement and recorded worker fit-hours describe a consolidation and
training-workload trade-off, not a demonstrated production benefit.

All conclusions remain bounded by source-specific eligible populations,
dependent temporal blocks, the existing seeds and the implemented model
families. Relative formulation gains remain distinct from absolute
persistence skill (Section~\ref{sec:results-persistence}). The post-hoc evidence
narrows generality and mechanism claims; it does not alter the original
scientific endpoints or extend applicability to excluded VMs and untested
shared architectures.

\section{Discussion}
\label{sec:discussion}

\subsection{When Cross-Resource Context Helps}
\label{sec:discussion-context}

Section~\ref{sec:results-rq1} motivates validating contextual inputs per
target, environment and time rather than assuming that additional telemetry
helps. The practical question is whether other histories improve the chosen
target in its operating context, not which resource should always receive them.

MMRP's univariate/multivariate single-target results likewise do not uniformly
favour multivariate inputs \citep{Tripathi2025MMRP}. This is a meaningful
precedent for the contextual comparison, although its reported feature and
selection details do not establish equivalence to our matched protocol.
Resource-forecasting and transfer experiments also make the evaluation
environment an explicit design consideration \citep{Rossi2025CloudUncertainty}.

TimeTrack's contemporaneous associations \citep{Meliani2025TimeTrack}
motivate candidate inputs but do not establish incremental predictive
information conditional on the target's own history. The comparisons neither
establish causal resource coupling nor explain observed directions through
intrinsic memory behaviour, throughput noise or disk-write sparsity; these
explanations require additional measurements.

\subsection{Shared Representation and Negative Transfer}
\label{sec:discussion-sharing}

Without independent all-input S2 models, the complete M1-versus-S1 comparison
changes input scope and representation sharing together. S2 is therefore an
important intermediate comparator, making contextual-input outcomes separately
observable from the implemented sharing transition under the matched
operational contrasts (Section~\ref{sec:results-rq2}).

Multi-task learning permits useful inductive transfer
\citep{Caruana1997Multitask}, but task-grouping research evaluates which
combinations benefit \citep{Standley2020TaskGrouping,Fifty2021TaskGroupings}.
Forecasting studies already examine negative transfer and selective sharing
\citep{Nikentari2022NARMAXLSTM,Nochumsohn2025MTLinear}; our contribution is a
scoped runtime-resource comparison, not discovery of the phenomenon.

The observed M1 degradation is operationally consistent with adverse transfer,
but the experiment does not identify the mechanism. Aggregate capacity,
model count, joint optimisation and common versus target-specific context
selection also differ between M1 and S2. These are matched operational
contrasts, not a parameter-matched causal test of sharing alone.
Gradient conflict, representation
interference, capacity competition and task covariance were not measured.
Multi-objective and gradient-aware approaches offer alternative training
designs \citep{Sener2018MultiObjective,Yu2020GradientSurgery}; their existence
does not diagnose our models. M1 is deliberately a conventional shared-LSTM
formulation. Its repeated adverse outcomes justify testing a common
representation before adopting it, not rejecting multi-task forecasting or
asserting that separate models are always preferable.

\subsection{Implications for Runtime and Cloud Software Systems}
\label{sec:discussion-systems}

Monitoring and forecasting inform software adaptation and resource management
\citep{Taherizadeh2018Monitoring,Bauer2020SelfAware}. Here, telemetry
selection and cross-target modelling should be separate decisions: additional
channels require validation, and a target benefiting from context may still
favour independently fitted prediction.

Evaluate local, contextual-independent and shared formulations separately
rather than treating ``multivariate'' as one category. Assess intended
resources and horizons under future-only validation, retaining environment-
and time-specific results. Relative gains do not establish absolute skill
(Section~\ref{sec:results-persistence}). This process neither prescribes one
architecture for every resource nor replaces system-specific validation.

Forecast-design evidence remains distinct from operational effectiveness.
Autoscaling research places prediction within a wider control loop
\citep{Qu2018Autoscaling}. We deployed no autoscaler, scheduler or controller
and measured no SLA improvement, monetary saving, energy reduction,
capacity-planning benefit or incident prevention. Those outcomes require
separate evaluation of forecasts together with decision policies and their
operating constraints.

\subsection{Accuracy-Efficiency Trade-Off}
\label{sec:discussion-efficiency}

Table~\ref{tab:efficiency} shows substantial consolidation: in the five-horizon
external instantiation, M1 uses approximately $74.5\%$ fewer parameters and
$75\%$ fewer models than S2, alongside lower observed Materna fitting workload.
Section~\ref{sec:results-efficiency} nevertheless records co-occurring accuracy
losses. Consolidation should therefore be evaluated as a trade-off, not a free
accuracy gain. Joint consideration of predictive quality and computational
requirements is established in task grouping and software-service forecasting
\citep{Standley2020TaskGrouping,Syu2017DynamicQoS}.

Counts and fit-hours characterise this implementation and workload, not
inference latency, deployed RAM consumption, energy or monetary cost.
Parameter reductions need not yield proportionate savings. Choosing a compact
model requires resource-specific error tolerances and direct measurements on
the intended platform; model count alone cannot establish a production saving.

\subsection{Applicability and Generalisation Limits}
\label{sec:discussion-applicability}

These implications concern four resource semantics and a conventional
two-layer LSTM, within particular runtime/cloud contexts. The external
populations are eligible dynamic multi-resource VMs, not unrestricted archive
samples. Bitbrains and Materna require nondegenerate training IQR across all
required scale-fit partitions; only 74 of Materna's 398 structurally common
VMs pass strict all-trace scale eligibility. The executed cohorts contain 35
selected VMs per Bitbrains environment and 35 shared Materna identifiers.
Section~\ref{sec:population-restrictions} explains why these findings do not
automatically extend to excluded idle or resource-constant VMs. Materna
T1--T3 remain dependent temporal blocks, not three independent datasets.

The results motivate testing selective resource grouping, adaptive sharing and
hybrid shared/task-specific representations. Task-grouping and gradient-aware
methods provide concrete candidates
\citep{Fifty2021TaskGroupings,Yu2020GradientSurgery,Nochumsohn2025MTLinear}.
Global forecasting research and runtime foundation-model comparisons also
motivate testing these alternatives
\citep{MonteroManso2021GlobalForecasting,DiMenna2026RuntimeForecasting}.
None of these alternatives, including
mixture-of-experts methods, was tested here, and none can be assumed to repair
M1's degradation. Future comparisons should retain independent all-input
controls and resource-specific assessment while measuring the mechanisms
they seek to explain.

\section{Threats to Validity}
\label{sec:threats}

\subsection{Construct Validity}
\label{sec:threats-construct}

The three arms are not a complete $2\times2$ factorial. S1/S2 differ in input
dimensionality and parameter count; S2/M1 differ in model count, aggregate
capacity, optimisation and target-specific versus common context selection.
The latter is therefore an operational shared-representation contrast, not a
parameter-matched causal intervention
(Section~\ref{sec:decomposition}). Its interpretation includes these
implementation differences, rather than attributing every change to shared
parameters alone.

CPU utilisation, throughput, memory usage and disk-write activity are semantic
mappings across datasets. Source definitions and native units differ, so the
mapping does not establish perfect measurement equivalence
(Section~\ref{sec:preprocessing}). Primary comparisons use MacroMAE-based
gains; RMSE, $R^2$, NMAE and persistence supply complementary diagnostics.
Metric choice may change apparent rankings \citep{Hyndman2006Accuracy}.
Accordingly, the reported accuracy boundary concerns the specified measures,
not every possible forecasting objective.

\subsection{Internal Validity}
\label{sec:threats-internal}

Temporal evaluation design affects forecasting performance estimation
\citep{Cerqueira2020Evaluation,Hyndman2021Forecasting}.
Section~\ref{sec:leakage-control} specifies chronological origins,
gap-bounded legal segments, expanding outer folds and maximum-horizon purging.
Inner validation remains inside outer-training support, scalers use training
data only, and outer tests do not select contexts or epochs. Matched support
across S1/S2/M1 reduces differences attributable to evaluation samples.
External and temporal criteria were fixed prospectively. These controls cannot
exclude every implementation or data-provenance error.

TimeTrack's large gap is handled through segmentation: windows and targets do
not cross it, but available contiguous support is reduced.
Three seeds limit dependence on one initialisation, while residual neural
optimisation stochasticity remains. Selected contexts and epochs may interact
with formulation. The models are not capacity matched, hyperparameter search
is intentionally limited, and the plain LSTM may favour or disfavour particular
sharing patterns (Sections~\ref{sec:architecture-training}
and~\ref{sec:model-selection}).

\subsection{External Validity}
\label{sec:threats-external}

TimeTrack, Bitbrains fastStorage and Rnd, and Materna T1--T3 represent
particular runtime/cloud contexts, not all workloads. Scale eligibility
retains 596 of 916 structurally eligible fastStorage VMs, 272 of 416 Rnd VMs,
and 74 of 398 structurally common Materna VMs. The strict rule requires
nondegenerate training IQR across every required resource/scaler partition,
including every Materna block. Only 35 VMs are selected per external
environment or cohort (Section~\ref{sec:population-restrictions}).
Deterministic hash selection is reproducible but is not probability sampling.
Conclusions therefore concern eligible dynamic multi-resource VM populations,
not every archived VM, particularly excluded idle or resource-constant cases.

Materna's consecutive blocks share the source and selected VM population,
providing repeated-within-source evidence rather than independent datasets
(Section~\ref{sec:replication-design}).

Only the controlled two-layer LSTM formulations are evaluated, including the
conventional shared M1 design. This does not establish
behaviour for Transformers, mixture-of-experts, selective task grouping,
gradient surgery, adaptive sharing or foundation forecasters. Only four
resource semantics are evaluated. Findings do not extend directly to latency,
energy, queue lengths, request rates, application QoS or SLA outcomes.

\subsection{Conclusion Validity}
\label{sec:threats-conclusion}

MAEs are averaged over seeds within matched cells
before equal weighting of VM--fold cells and horizons
(Section~\ref{sec:evaluation}). External uncertainty uses paired VM-cluster
bootstrap intervals. These intervals condition on the evaluated cohorts and
protocol; they do not establish universal population inference.
Conclusions consider magnitude and frozen consistency criteria rather than
one aggregate number. Nevertheless, multiple resource/horizon comparisons
create many opportunities for variation. Intervals and classifications are
not presented as classical null-hypothesis significance claims.

Persistence distinguishes gains over another learned formulation from absolute
forecasting skill (Section~\ref{sec:results-persistence}). Bitbrains and
Materna syntheses are descriptive equal-environment and equal-temporal-block
summaries, respectively. The six-block disk-write mean in
Section~\ref{sec:cross-dataset-synthesis} is likewise descriptive only, not an
effect estimated from independent datasets. No cross-dataset inferential
meta-analysis is performed. This is formulation- and population-bounded
evidence, not a universal statistical verdict on resource sharing.

\section{Conclusion}
\label{sec:conclusion}

Multi-metric runtime forecasting involves two distinct decisions: whether a
predictor receives other resources' histories, and whether resource targets
share a learned predictive representation. S2 versus S1 evaluates contextual
information, M1 versus S2 implemented representation sharing, and M1 versus S1 the complete
joint-system change. These operational contrasts were examined through
TimeTrack development, Bitbrains external environments and Materna dependent
temporal replication.

Cross-resource context was selective, varying with resource, environment and
time. Memory showed the clearest repeated contextual benefit, but not in every
block. For the evaluated M1 shared-LSTM formulation, throughput, memory and disk write had
adverse M1-versus-S2 effects across all six evidence blocks; CPU contained
limited exceptions. The complete M1 system generally did not outperform local
S1 forecasting, with Materna T2 CPU the exception. These accuracy losses
coexisted with substantial consolidation: in the five-horizon external instantiation, M1 used
approximately $74.5\%$ fewer parameters and $75\%$ fewer models than S2. The
result is an accuracy--efficiency boundary for the evaluated implementation,
not a general rejection of shared modelling.

Cross-resource histories should be validated per target, and shared
representations assessed separately against independently
trained contextual models. The scope is the four evaluated resource semantics
and conventional shared-LSTM formulation, with external conclusions restricted
to eligible dynamic VM populations. The evidence does not establish
performance for all cloud systems, alternative multi-task architectures or
production controllers. Future work should test whether selective or adaptive
sharing can retain consolidation benefits while avoiding the adverse transfer
observed here.

\section*{Data availability}

TimeTrack is publicly available on
\href{https://www.kaggle.com/datasets/abdelghanimeliani/open-air-interface-cicd-timeseries-dataset}{Kaggle}
(Apache 2.0 in its dataset metadata) and in the publisher's
\href{https://doi.org/10.5281/zenodo.15442251}{Zenodo archive
(doi:10.5281/zenodo.15442251; CC BY 4.0)}; its description paper is
\citet{Meliani2025TimeTrack}. These are source-specific licence statements.
Bitbrains fastStorage and Rnd are publicly available as
\href{https://atlarge-research.com/gwa-t-12/}{GWA-T-12}
\citep{Shen2015Bitbrains}; Materna T1--T3 are publicly available as
\href{https://atlarge-research.com/gwa-t-13/}{GWA-T-13}
\citep{Kohne2014FederatedCloudSim,Kohne2016SLAScheduling}.
We acknowledge Bitbrains IT Services Inc. and Materna GmbH Information
\& Communications, respectively, for providing these traces through the
Grid Workloads Archive. Their separate source pages specify acknowledgement
and citation conditions, rather than a standard open-data licence. Readers
should obtain third-party raw data from these original sources and consult
each source's terms before redistribution.

Source code, protocols, derived analysis outputs and SHA-256 inventories are
publicly accessible in the
\href{https://github.com/KADOUMAAbdelhak/TimeTrack-Forecasting-Framework/tree/fee20c316040c57b69ba12a193530c955b539771}{project repository at scientific checkpoint \texttt{fee20c3}}.
The public checkpoint does not include the full fitted-model or per-fit
prediction collections; the checksum inventories do not themselves provide
those files. No repository-wide reuse licence is asserted here.

\printcredits

\section*{Funding}

This research did not receive any specific grant from funding agencies in
the public, commercial, or not-for-profit sectors.


\bibliographystyle{cas-model2-names}
\bibliography{references}

\end{document}
